%%%%%%%%%%%%%%%%%%%%%%%%%%%%%%%%%%%%%%%%%%%%%%%%%%%
%% Exercices & Solutions — Informatique Quantique
%% Stage LIA / CERI — Université d'Avignon
%% Stagiaire : Patrice SEBASTIANO
%% Encadrant : Serigne GUEYE
%%%%%%%%%%%%%%%%%%%%%%%%%%%%%%%%%%%%%%%%%%%%%%%%%%%

% \documentclass[light]{ceri/sty/rapport}   % compilation rapide
% \documentclass[handout]{ceri/sty/rapport}  % impression
% \documentclass[full]{ceri/sty/rapport}     % avec listes des figures/tableaux
\documentclass{ceri/sty/rapport}


%%%%%%%%%%%%%%%%%%%%%%%%%%%%%%%%%%%%%%%%%%%%%%%%%%%
%% Informations générales
%%%%%%%%%%%%%%%%%%%%%%%%%%%%%%%%%%%%%%%%%%%%%%%%%%%
\major{Master d'Informatique}
\specialization{Informatique Quantique}
\course{Stage de Recherche — LIA Avignon}

\title{Exercices \& Solutions}
\subtitle{Informatique Quantique}

\author{
	Patrice SEBASTIANO
}

\advisor[Encadrant]{
	Serigne GUEYE
}


%%%%%%%%%%%%%%%%%%%%%%%%%%%%%%%%%%%%%%%%%%%%%%%%%%%
%% Packages supplémentaires
%%%%%%%%%%%%%%%%%%%%%%%%%%%%%%%%%%%%%%%%%%%%%%%%%%%
\usepackage{ifthen}
\usepackage{amsmath, amssymb}
\usepackage{mathtools}
\usepackage{braket}
\usepackage{booktabs}
\usepackage{changepage}
\usepackage{quantikz}
\usepackage{tcolorbox}
\tcbuselibrary{breakable, skins}

% Couleurs personnalisées
\definecolor{exo-bg}{RGB}{230, 240, 255}
\definecolor{exo-border}{RGB}{50, 100, 200}
\definecolor{sol-bg}{RGB}{230, 255, 235}
\definecolor{sol-border}{RGB}{40, 160, 60}

% Environnement exercice — argument optionnel = référence page du livre
\newcounter{exercice}[section]
\renewcommand{\theexercice}{\thesection.\arabic{exercice}}

\newtcolorbox[use counter=exercice]{exercice}[1][]{%
	colback=exo-bg,
	colframe=exo-border,
	fonttitle=\bfseries,
	title={Exercice~\theexercice \ifthenelse{\equal{#1}{}}{}{~—~#1}},
	breakable,
	before upper={\parindent=1em}
}

% Environnement solution — argument optionnel transmis à tcolorbox (ex: élargissement)
\newtcolorbox{solution}[1][]{%
	colback=sol-bg,
	colframe=sol-border,
	fonttitle=\bfseries,
	title={Solution},
	breakable,
	before upper={\parindent=1em},
	#1
}


%%%%%%%%%%%%%%%%%%%%%%%%%%%%%%%%%%%%%%%%%%%%%%%%%%%
%% Début du document
%%%%%%%%%%%%%%%%%%%%%%%%%%%%%%%%%%%%%%%%%%%%%%%%%%%
\begin{document}

\maketitle
\sloppy


%%%%%%%%%%%%%%%%%%%%%%%%%%%%%%%%%%%%%%%%%%%%%%%%%%%
\section{Exercices du Chapitre 1}
\label{sec:chap1}
%%%%%%%%%%%%%%%%%%%%%%%%%%%%%%%%%%%%%%%%%%%%%%%%%%%

\begin{exercice}[p.~11]
What is the probability of measuring $0$ if the state of a qubit is
$\sqrt{1/2}\,\ket{0} + \sqrt{1/2}\,\ket{1}$?
And the probability of measuring $1$?
What if the state of the qubit is $\sqrt{1/3}\,\ket{0} + \sqrt{2/3}\,\ket{1}$?
And if it is $\sqrt{1/2}\,\ket{0} - \sqrt{1/2}\,\ket{1}$?
\end{exercice}

\begin{solution}
For a state $\alpha\ket{0}+\beta\ket{1}$, the Born rule gives $P(0)=|\alpha|^2$ and
$P(1)=|\beta|^2$.

\medskip\noindent
\textbf{State $\sqrt{1/2}\,\ket{0}+\sqrt{1/2}\,\ket{1}$:}
$P(0)=\tfrac{1}{2}$, $P(1)=\tfrac{1}{2}$.

\medskip\noindent
\textbf{State $\sqrt{1/3}\,\ket{0}+\sqrt{2/3}\,\ket{1}$:}
$P(0)=\tfrac{1}{3}$, $P(1)=\tfrac{2}{3}$.

\medskip\noindent
\textbf{State $\sqrt{1/2}\,\ket{0}-\sqrt{1/2}\,\ket{1}$:}
$P(0)=\bigl(\sqrt{1/2}\bigr)^2=\tfrac{1}{2}$, $P(1)=\bigl(-\sqrt{1/2}\bigr)^2=\tfrac{1}{2}$.
The sign of the amplitude does not affect measurement probabilities.
\end{solution}

\begin{exercice}[p.~13]
What is the inner product of $\sqrt{1/2}\,\ket{0} + \sqrt{1/2}\,\ket{1}$ and
$\sqrt{1/3}\,\ket{0} + \sqrt{2/3}\,\ket{1}$?
And the inner product of $\sqrt{1/2}\,\ket{0} + \sqrt{1/2}\,\ket{1}$ and
$\sqrt{1/2}\,\ket{0} - \sqrt{1/2}\,\ket{1}$?
\end{exercice}

\begin{solution}
Recall $\braket{0|0}=\braket{1|1}=1$ and $\braket{0|1}=\braket{1|0}=0$.

\medskip\noindent
\textbf{First pair:}
\[
  \Bigl(\sqrt{\tfrac{1}{2}}\bra{0}+\sqrt{\tfrac{1}{2}}\bra{1}\Bigr)
  \Bigl(\sqrt{\tfrac{1}{3}}\ket{0}+\sqrt{\tfrac{2}{3}}\ket{1}\Bigr)
  = \sqrt{\tfrac{1}{2}}\cdot\sqrt{\tfrac{1}{3}} + \sqrt{\tfrac{1}{2}}\cdot\sqrt{\tfrac{2}{3}}
  = \frac{1}{\sqrt{6}} + \frac{\sqrt{2}}{\sqrt{6}}
  = \frac{1+\sqrt{2}}{\sqrt{6}}.
\]

\medskip\noindent
\textbf{Second pair:}
\[
  \Bigl(\sqrt{\tfrac{1}{2}}\bra{0}+\sqrt{\tfrac{1}{2}}\bra{1}\Bigr)
  \Bigl(\sqrt{\tfrac{1}{2}}\ket{0}-\sqrt{\tfrac{1}{2}}\ket{1}\Bigr)
  = \frac{1}{2} - \frac{1}{2} = 0.
\]
The two states $\ket{+}$ and $\ket{-}$ are orthogonal.
\end{solution}

\begin{exercice}[p.~15]
Check that the gate $X$ matrix is, indeed, unitary.
What is the inverse of $X$?
What is the action of $X$ on a general qubit in a state of the form
$a\ket{0} + b\ket{1}$?
\end{exercice}

\begin{solution}
The $X$ gate has matrix $X=\begin{pmatrix}0&1\\1&0\end{pmatrix}$.
Since $X$ is real and symmetric, $X^\dagger = X^\top = X$, so
\[
  X^\dagger X = X^2 =
  \begin{pmatrix}0&1\\1&0\end{pmatrix}
  \begin{pmatrix}0&1\\1&0\end{pmatrix}
  = \begin{pmatrix}1&0\\0&1\end{pmatrix} = I.
\]
Hence $X$ is unitary.  Since $X^2=I$, the inverse is $X^{-1}=X$ (the gate is its own inverse).

The action on a general qubit:
\[
  X\bigl(a\ket{0}+b\ket{1}\bigr) = aX\ket{0}+bX\ket{1} = a\ket{1}+b\ket{0}.
\]
$X$ swaps the amplitudes of $\ket{0}$ and $\ket{1}$, acting as a classical NOT gate.
\end{solution}

\begin{exercice}[p.~16]
Check that the gate $H$ matrix is, indeed, unitary.
What is the action of $H$ on $\ket{+}$ and $\ket{-}$?
What is the action of $X$ on $\ket{+}$ and $\ket{-}$?
\end{exercice}

\begin{solution}
The Hadamard matrix is $H=\tfrac{1}{\sqrt{2}}\begin{pmatrix}1&1\\1&-1\end{pmatrix}$.
It is real and symmetric, so $H^\dagger=H$, and
\[
  H^2 = \frac{1}{2}\begin{pmatrix}1&1\\1&-1\end{pmatrix}^2
      = \frac{1}{2}\begin{pmatrix}2&0\\0&2\end{pmatrix} = I,
\]
hence $H^\dagger H = I$: $H$ is unitary.

\medskip\noindent
\textbf{Action of $H$:}
Using $H\ket{0}=\ket{+}$ and $H\ket{1}=\ket{-}$,
\[
  H\ket{+} = H\cdot\frac{\ket{0}+\ket{1}}{\sqrt{2}}
           = \frac{H\ket{0}+H\ket{1}}{\sqrt{2}}
           = \frac{\ket{+}+\ket{-}}{\sqrt{2}}
           = \frac{(\ket{0}+\ket{1})+(\ket{0}-\ket{1})}{2} = \ket{0}.
\]
\[
  H\ket{-} = \frac{H\ket{0}-H\ket{1}}{\sqrt{2}} = \frac{\ket{+}-\ket{-}}{\sqrt{2}}
           = \frac{2\ket{1}/\sqrt{2}}{\sqrt{2}} = \ket{1}.
\]
$H$ swaps the two bases: $\ket{+}\leftrightarrow\ket{0}$, $\ket{-}\leftrightarrow\ket{1}$.

\medskip\noindent
\textbf{Action of $X$:}
\[
  X\ket{+} = X\cdot\frac{\ket{0}+\ket{1}}{\sqrt{2}} = \frac{\ket{1}+\ket{0}}{\sqrt{2}} = \ket{+}.
\]
\[
  X\ket{-} = X\cdot\frac{\ket{0}-\ket{1}}{\sqrt{2}} = \frac{\ket{1}-\ket{0}}{\sqrt{2}} = -\ket{-}.
\]
$\ket{+}$ is an eigenvector of $X$ with eigenvalue $+1$; $\ket{-}$ with eigenvalue $-1$.
\end{solution}

\begin{exercice}[p.~17]
Check that $Z\ket{0} = \ket{0}$ and that $Z\ket{1} = -\ket{1}$ in two different
ways. First, use Dirac notation and the actions of $H$ and $X$ (remember that we
have defined $Z$ as $HXH$). Then, derive the same result by performing the matrix
multiplication
\[
  \begin{pmatrix}
    \dfrac{1}{\sqrt{2}} & \dfrac{1}{\sqrt{2}} \\[6pt]
    \dfrac{1}{\sqrt{2}} & -\dfrac{1}{\sqrt{2}}
  \end{pmatrix}
  \begin{pmatrix} 0 & 1 \\ 1 & 0 \end{pmatrix}
  \begin{pmatrix}
    \dfrac{1}{\sqrt{2}} & \dfrac{1}{\sqrt{2}} \\[6pt]
    \dfrac{1}{\sqrt{2}} & -\dfrac{1}{\sqrt{2}}
  \end{pmatrix}.
\]
\end{exercice}

\begin{solution}
\textbf{Method 1 — Dirac notation.}
Using $Z=HXH$, $H\ket{0}=\ket{+}$, $H\ket{1}=\ket{-}$, $X\ket{+}=\ket{+}$, $X\ket{-}=-\ket{-}$:
\[
  Z\ket{0} = HXH\ket{0} = HX\ket{+} = H\ket{+} = \ket{0}.\qquad\checkmark
\]
\[
  Z\ket{1} = HXH\ket{1} = HX\ket{-} = H(-\ket{-}) = -H\ket{-} = -\ket{1}.\qquad\checkmark
\]

\medskip\noindent
\textbf{Method 2 — Matrix multiplication.}
With $H=\tfrac{1}{\sqrt{2}}\begin{pmatrix}1&1\\1&-1\end{pmatrix}$, first compute $XH$:
\[
  XH = \frac{1}{\sqrt{2}}\begin{pmatrix}0&1\\1&0\end{pmatrix}\begin{pmatrix}1&1\\1&-1\end{pmatrix}
     = \frac{1}{\sqrt{2}}\begin{pmatrix}1&-1\\1&1\end{pmatrix}.
\]
Then $H(XH)$:
\[
  HXH = \frac{1}{2}\begin{pmatrix}1&1\\1&-1\end{pmatrix}\begin{pmatrix}1&-1\\1&1\end{pmatrix}
      = \frac{1}{2}\begin{pmatrix}2&0\\0&-2\end{pmatrix}
      = \begin{pmatrix}1&0\\0&-1\end{pmatrix} = Z.\qquad\checkmark
\]
\end{solution}

\begin{exercice}[p.~18]
Check that $T^2 = S$. Then, use the most beautiful formula ever
(i.e., Euler's identity $e^{i\pi}+1=0$) to check that $S^2 = Z$.
Check also that $S$ and $T$ are unitary.
Express $S^\dagger$ and $T^\dagger$ as powers of $S$ and $T$.
\end{exercice}

\begin{solution}
Recall $T=\begin{pmatrix}1&0\\0&e^{i\pi/4}\end{pmatrix}$ and $S=\begin{pmatrix}1&0\\0&i\end{pmatrix}$.

\medskip\noindent
\textbf{$T^2=S$:}
\[
  T^2 = \begin{pmatrix}1&0\\0&e^{i\pi/4}\end{pmatrix}^2
      = \begin{pmatrix}1&0\\0&e^{i\pi/2}\end{pmatrix}
      = \begin{pmatrix}1&0\\0&i\end{pmatrix} = S.\qquad\checkmark
\]

\medskip\noindent
\textbf{$S^2=Z$:}
\[
  S^2 = \begin{pmatrix}1&0\\0&i\end{pmatrix}^2
      = \begin{pmatrix}1&0\\0&i^2\end{pmatrix}.
\]
Since $i^2=-1$, we get
$S^2=\begin{pmatrix}1&0\\0&-1\end{pmatrix}=Z$.\quad$\checkmark$

\medskip\noindent
\textbf{Unitarity:}
Both $S$ and $T$ are diagonal with entries of modulus~$1$, so their adjoints are
$S^\dagger=\begin{pmatrix}1&0\\0&-i\end{pmatrix}$ and $T^\dagger=\begin{pmatrix}1&0\\0&e^{-i\pi/4}\end{pmatrix}$.
One immediately verifies $S^\dagger S = T^\dagger T = I$.

\medskip\noindent
\textbf{Powers:}
Since $S^4=(S^2)^2=Z^2=I$, we have $S^{-1}=S^3$, hence $S^\dagger = S^3$.\\
Since $T^8=(T^4)^2=(S^2)^2=I$, we have $T^{-1}=T^7$, hence $T^\dagger = T^7$.
\end{solution}

\begin{exercice}[p.~21]
Check these equivalences by substituting the angles in the definitions of
$R_X$, $R_Y$, and $R_Z$:
\[
  R_X(\pi) = X, \qquad R_Y(\pi) = Y, \qquad R_Z(\pi) = Z,
\]
\[
  R_Z\!\left(\frac{\pi}{2}\right) = S, \qquad R_Z\!\left(\frac{\pi}{4}\right) = T.
\]
\end{exercice}

\begin{solution}
Using the standard definitions
$R_X(\theta)=\begin{pmatrix}\cos\frac{\theta}{2}&-i\sin\frac{\theta}{2}\\-i\sin\frac{\theta}{2}&\cos\frac{\theta}{2}\end{pmatrix}$,
$R_Y(\theta)=\begin{pmatrix}\cos\frac{\theta}{2}&-\sin\frac{\theta}{2}\\\sin\frac{\theta}{2}&\cos\frac{\theta}{2}\end{pmatrix}$,
$R_Z(\theta)=\begin{pmatrix}e^{-i\theta/2}&0\\0&e^{i\theta/2}\end{pmatrix}$.

\[
  R_X(\pi)=\begin{pmatrix}0&-i\\-i&0\end{pmatrix}=-i\begin{pmatrix}0&1\\1&0\end{pmatrix}=-iX,\qquad
  R_Y(\pi)=\begin{pmatrix}0&-1\\1&0\end{pmatrix}=-iY,
\]
\[
  R_Z(\pi)=\begin{pmatrix}-i&0\\0&i\end{pmatrix}=-i\begin{pmatrix}1&0\\0&-1\end{pmatrix}=-iZ.
\]
Each rotation equals the corresponding Pauli gate \emph{up to the global phase} $-i$, which is unobservable.
\[
  R_Z\!\left(\tfrac{\pi}{2}\right)
    =\begin{pmatrix}e^{-i\pi/4}&0\\0&e^{i\pi/4}\end{pmatrix}
    =e^{-i\pi/4}\begin{pmatrix}1&0\\0&e^{i\pi/2}\end{pmatrix}
    =e^{-i\pi/4}S,
\]
\[
  R_Z\!\left(\tfrac{\pi}{4}\right)
    =e^{-i\pi/8}\begin{pmatrix}1&0\\0&e^{i\pi/4}\end{pmatrix}
    =e^{-i\pi/8}T.
\]
Again, equality with $S$ and $T$ holds up to a global phase.
\end{solution}

\begin{exercice}[p.~22]
Show that $U(\theta,\varphi,\lambda)$ is unitary.
Check that $R_X(\theta) = U(\theta,-\pi/2,\pi/2)$, that
$R_Y(\theta) = U(\theta,0,0)$ and that, up to a global phase,
$R_Z(\theta) = U(0,0,\theta)$.
\end{exercice}

\begin{solution}
The universal single-qubit gate is defined as
\[
  U(\theta,\varphi,\lambda)
  =\begin{pmatrix}
      \cos\dfrac{\theta}{2} & -e^{i\lambda}\sin\dfrac{\theta}{2}\\[8pt]
      e^{i\varphi}\sin\dfrac{\theta}{2} & e^{i(\varphi+\lambda)}\cos\dfrac{\theta}{2}
   \end{pmatrix}.
\]

\medskip\noindent
\textbf{Unitarity.}
Set $c=\cos\tfrac{\theta}{2}$ and $s=\sin\tfrac{\theta}{2}$ for brevity.
The conjugate transpose is
\[
  U^\dagger
  =\begin{pmatrix}
      c           & e^{-i\varphi}s         \\[4pt]
      -e^{-i\lambda}s & e^{-i(\varphi+\lambda)}c
   \end{pmatrix}.
\]
We compute $U^\dagger U$ entry by entry.
\begin{align*}
  (U^\dagger U)_{11}
    &= c\cdot c + e^{-i\varphi}s\cdot e^{i\varphi}s
     = c^2+s^2 = 1.\\[4pt]
  (U^\dagger U)_{12}
    &= c\cdot(-e^{i\lambda}s) + e^{-i\varphi}s\cdot e^{i(\varphi+\lambda)}c
     = -e^{i\lambda}cs + e^{i\lambda}cs = 0.\\[4pt]
  (U^\dagger U)_{21}
    &= -e^{-i\lambda}s\cdot c + e^{-i(\varphi+\lambda)}c\cdot e^{i\varphi}s
     = -e^{-i\lambda}cs + e^{-i\lambda}cs = 0.\\[4pt]
  (U^\dagger U)_{22}
    &= (-e^{-i\lambda}s)(-e^{i\lambda}s) + e^{-i(\varphi+\lambda)}c\cdot e^{i(\varphi+\lambda)}c
     = s^2+c^2 = 1.
\end{align*}
Hence $U^\dagger U = I$, so $U$ is unitary.\quad$\checkmark$

\medskip\noindent
\textbf{$R_X(\theta) = U(\theta,\,-\pi/2,\,\pi/2)$.}
Substitute $\varphi=-\pi/2$ and $\lambda=\pi/2$:
\begin{align*}
  -e^{i\lambda}s &= -e^{i\pi/2}s = -is,\\
  e^{i\varphi}s  &= e^{-i\pi/2}s = -is,\\
  e^{i(\varphi+\lambda)}c &= e^{i(-\pi/2+\pi/2)}c = e^{0}c = c.
\end{align*}
Therefore
\[
  U(\theta,-\tfrac{\pi}{2},\tfrac{\pi}{2})
  = \begin{pmatrix}c & -is\\-is & c\end{pmatrix}
  = R_X(\theta).\quad\checkmark
\]

\medskip\noindent
\textbf{$R_Y(\theta) = U(\theta,\,0,\,0)$.}
Substitute $\varphi=0$ and $\lambda=0$:
\[
  U(\theta,0,0)
  = \begin{pmatrix}c & -s\\s & c\end{pmatrix}
  = R_Y(\theta).\quad\checkmark
\]

\medskip\noindent
\textbf{$R_Z(\theta) = U(0,\,0,\,\theta)$ up to a global phase.}
Substitute $\theta=0$ (so $c=1$, $s=0$) and $\lambda=\theta$:
\[
  U(0,0,\theta)
  = \begin{pmatrix}1 & 0\\0 & e^{i\theta}\end{pmatrix}.
\]
Factor out $e^{i\theta/2}$:
\[
  U(0,0,\theta)
  = e^{i\theta/2}
    \begin{pmatrix}e^{-i\theta/2} & 0\\0 & e^{i\theta/2}\end{pmatrix}
  = e^{i\theta/2}\,R_Z(\theta).
\]
The two gates differ only by the global phase $e^{i\theta/2}$, which is
physically unobservable.\quad$\checkmark$
\end{solution}

\begin{exercice}[p.~23]
Check that, with the preceding circuit, the state before measurement is
$\sqrt{p}\,\ket{0} + \sqrt{1-p}\,\ket{1}$ and, hence, the probability of
measuring $0$ is $p$ and that of measuring $1$ is $1-p$.
\end{exercice}

\begin{solution}
The circuit is:
\[
\begin{quantikz}
  \lstick{$\ket{0}$} & \gate{R_Y(\theta)} & \meter{}
\end{quantikz}
\]

\medskip\noindent
\textbf{Choice of $\theta$.}
We set $\theta = 2\arccos\!\left(\sqrt{p}\right)$, so that
\[
  \frac{\theta}{2} = \arccos\!\left(\sqrt{p}\right)
  \quad\Longrightarrow\quad
  \cos\frac{\theta}{2} = \sqrt{p}.
\]
From the Pythagorean identity $\sin^2\!\frac{\theta}{2}+\cos^2\!\frac{\theta}{2}=1$:
\[
  \sin\frac{\theta}{2}
  = \sqrt{1-\cos^2\!\frac{\theta}{2}}
  = \sqrt{1-p}.
\]

\medskip\noindent
\textbf{Explicit computation of $R_Y(\theta)\ket{0}$.}
\[
  R_Y(\theta)\ket{0}
  =\begin{pmatrix}\cos\dfrac{\theta}{2} & -\sin\dfrac{\theta}{2}\\[8pt]
                  \sin\dfrac{\theta}{2} &  \cos\dfrac{\theta}{2}\end{pmatrix}
   \begin{pmatrix}1\\0\end{pmatrix}
  =\begin{pmatrix}\cos\dfrac{\theta}{2}\\[8pt]\sin\dfrac{\theta}{2}\end{pmatrix}
  =\begin{pmatrix}\sqrt{p}\\[4pt]\sqrt{1-p}\end{pmatrix}
  =\sqrt{p}\,\ket{0}+\sqrt{1-p}\,\ket{1}.
\]

\medskip\noindent
\textbf{Probabilities.}
By the Born rule applied to the state $\sqrt{p}\,\ket{0}+\sqrt{1-p}\,\ket{1}$:
\[
  P(\text{measure }0) = \left(\sqrt{p}\right)^2 = p,\qquad
  P(\text{measure }1) = \left(\sqrt{1-p}\right)^2 = 1-p.\quad\checkmark
\]
\end{solution}

\begin{exercice}[p.~26]
Derive the formulas for the probability of measuring $1$ on the first qubit
in a general two-qubit state and for the state of the system after the
measurement.
\end{exercice}

\begin{solution}
Write the general two-qubit state as
$\ket{\psi}=\alpha_{00}\ket{00}+\alpha_{01}\ket{01}+\alpha_{10}\ket{10}+\alpha_{11}\ket{11}$.
Grouping by the value of the first qubit:
\[
  \ket{\psi}=\ket{0}\!\otimes\!(\alpha_{00}\ket{0}+\alpha_{01}\ket{1})
            +\ket{1}\!\otimes\!(\alpha_{10}\ket{0}+\alpha_{11}\ket{1}).
\]
The probability of getting outcome $1$ on the first qubit is
\[
  P(q_1=1) = |\alpha_{10}|^2+|\alpha_{11}|^2.
\]
After the measurement the state collapses to the normalised projection onto the
first-qubit-$1$ subspace:
\[
  \ket{\psi'} = \frac{\alpha_{10}\ket{10}+\alpha_{11}\ket{11}}{\sqrt{|\alpha_{10}|^2+|\alpha_{11}|^2}}
              = \ket{1}\otimes\frac{\alpha_{10}\ket{0}+\alpha_{11}\ket{1}}{\sqrt{|\alpha_{10}|^2+|\alpha_{11}|^2}}.
\]
\end{solution}

\begin{exercice}[p.~27]
Check that, given any pair of unitary matrices $U_1$ and $U_2$, the inverse
of $U_1 \otimes U_2$ is $U_1^\dagger \otimes U_2^\dagger$ and that
$(U_1 \otimes U_2)^\dagger = U_1^\dagger \otimes U_2^\dagger$.
\end{exercice}

\begin{solution}
\textbf{Adjoint of $U_1\otimes U_2$.}
Write $U_k=\begin{pmatrix}a_k&b_k\\c_k&d_k\end{pmatrix}$ for $k=1,2$.
By definition of the Kronecker product:
\[
  U_1\otimes U_2
  =\begin{pmatrix}a_1 U_2 & b_1 U_2\\c_1 U_2 & d_1 U_2\end{pmatrix}.
\]
Taking the conjugate-transpose means transposing the block structure \emph{and}
conjugate-transposing each block:
\[
  (U_1\otimes U_2)^\dagger
  =\begin{pmatrix}\bar a_1 U_2^\dagger & \bar c_1 U_2^\dagger\\
                  \bar b_1 U_2^\dagger & \bar d_1 U_2^\dagger\end{pmatrix}.
\]
On the other hand,
$U_1^\dagger=\begin{pmatrix}\bar a_1 & \bar c_1\\\bar b_1 & \bar d_1\end{pmatrix}$,
so
\[
  U_1^\dagger\otimes U_2^\dagger
  =\begin{pmatrix}\bar a_1 U_2^\dagger & \bar c_1 U_2^\dagger\\
                  \bar b_1 U_2^\dagger & \bar d_1 U_2^\dagger\end{pmatrix}.
\]
The two expressions are identical, hence
$(U_1\otimes U_2)^\dagger = U_1^\dagger\otimes U_2^\dagger$.\quad$\checkmark$

\medskip\noindent
\textbf{Unitarity.}
Using the \emph{mixed-product property} $(A\otimes B)(C\otimes D)=(AC)\otimes(BD)$:
\[
  (U_1^\dagger\otimes U_2^\dagger)(U_1\otimes U_2)
  =(U_1^\dagger U_1)\otimes(U_2^\dagger U_2)
  =I\otimes I=I.\quad\checkmark
\]
\end{solution}

\begin{exercice}[p.~28]
Explicitly compute the matrices for the gates $X \otimes X$ and $H \otimes I$.
\end{exercice}

\begin{solution}
Using the block-matrix formula $A\otimes B = (a_{ij}B)_{ij}$, with basis ordering
$\ket{00},\ket{01},\ket{10},\ket{11}$:

\[
  X\otimes X =
  \begin{pmatrix}0\cdot X & 1\cdot X\\1\cdot X & 0\cdot X\end{pmatrix}
  =\begin{pmatrix}0&0&0&1\\0&0&1&0\\0&1&0&0\\1&0&0&0\end{pmatrix}.
\]

\[
  H\otimes I = \frac{1}{\sqrt{2}}
  \begin{pmatrix}1\cdot I & 1\cdot I\\1\cdot I & -1\cdot I\end{pmatrix}
  =\frac{1}{\sqrt{2}}\begin{pmatrix}1&0&1&0\\0&1&0&1\\1&0&-1&0\\0&1&0&-1\end{pmatrix}.
\]
\end{solution}

\begin{exercice}[p.~30]
Check these equivalences in two different ways: by computing the matrices of
the circuits and by obtaining the result of using them with qubits in states
$\ket{00}$, $\ket{01}$, $\ket{10}$, and $\ket{11}$.
\end{exercice}

\begin{solution}
The equivalence to verify is that three CNOTs implement a SWAP gate:
\[
  \mathrm{CNOT}_{1\to2}\;\mathrm{CNOT}_{2\to1}\;\mathrm{CNOT}_{1\to2}
  \;=\;\mathrm{SWAP}.
\]
In circuit form:
\[
\begin{quantikz}
  \qw & \ctrl{1}  & \targ{}    & \ctrl{1}  & \qw \\
  \qw & \targ{}   & \ctrl{-1}  & \targ{}   & \qw
\end{quantikz}
\quad=\quad
\begin{quantikz}
  \qw & \swap{1} & \qw \\
  \qw & \targX{} & \qw
\end{quantikz}
\]

\textbf{Method 1 — matrices.}

We take $\mathrm{CNOT}_{1\to2}$ as known:
\[
  M_1 = \mathrm{CNOT}_{1\to2}
      = \begin{pmatrix}1&0&0&0\\0&1&0&0\\0&0&0&1\\0&0&1&0\end{pmatrix}.
\]

For $\mathrm{CNOT}_{2\to1}$ (control $= q_2$, target $= q_1$) we start from a generic matrix
$M_2=\begin{pmatrix}a&b&c&d\\e&f&g&h\\i&j&k&l\\m&n&o&p\end{pmatrix}$
and impose its action on each basis state.

The gate acts as: flip $q_1$ iff $q_2=1$, i.e.
\[
  M_2\ket{00}=\ket{00},\quad
  M_2\ket{01}=\ket{11},\quad
  M_2\ket{10}=\ket{10},\quad
  M_2\ket{11}=\ket{01}.
\]

Each condition fixes one column of $M_2$ (a matrix column is the image of a basis vector):
\begin{align*}
  M_2\ket{00}=\ket{00} &\;\Rightarrow\; \text{col.\,1} = (1,0,0,0)^\top,\\
  M_2\ket{01}=\ket{11} &\;\Rightarrow\; \text{col.\,2} = (0,0,0,1)^\top,\\
  M_2\ket{10}=\ket{10} &\;\Rightarrow\; \text{col.\,3} = (0,0,1,0)^\top,\\
  M_2\ket{11}=\ket{01} &\;\Rightarrow\; \text{col.\,4} = (0,1,0,0)^\top.
\end{align*}

Hence all 16 unknowns are determined:
\[
  M_2 = \mathrm{CNOT}_{2\to1}
      = \begin{pmatrix}1&0&0&0\\0&0&0&1\\0&0&1&0\\0&1&0&0\end{pmatrix}.
\]

We now multiply the three matrices $M_1 M_2 M_1$ (rightmost applied first).
First, $M_2 M_1$:
\[
  M_2 M_1
  =\begin{pmatrix}1&0&0&0\\0&0&0&1\\0&0&1&0\\0&1&0&0\end{pmatrix}
   \begin{pmatrix}1&0&0&0\\0&1&0&0\\0&0&0&1\\0&0&1&0\end{pmatrix}
  =\begin{pmatrix}1&0&0&0\\0&0&1&0\\0&0&0&1\\0&1&0&0\end{pmatrix}.
\]
Then $M_1(M_2 M_1)$:
\[
  M_1 M_2 M_1
  =\begin{pmatrix}1&0&0&0\\0&1&0&0\\0&0&0&1\\0&0&1&0\end{pmatrix}
   \begin{pmatrix}1&0&0&0\\0&0&1&0\\0&0&0&1\\0&1&0&0\end{pmatrix}
  =\begin{pmatrix}1&0&0&0\\0&0&1&0\\0&1&0&0\\0&0&0&1\end{pmatrix}
  =\mathrm{SWAP}.\quad\checkmark
\]

\medskip\noindent
\textbf{Method 2 — basis states.}
We trace each basis state through the three gates
$\mathrm{CNOT}_{1\to2}\to\mathrm{CNOT}_{2\to1}\to\mathrm{CNOT}_{1\to2}$
and compare with $\mathrm{SWAP}\ket{q_1 q_2}=\ket{q_2 q_1}$:
\begin{align*}
\ket{00}&\xrightarrow{M_1}\ket{00}\xrightarrow{M_2}\ket{00}\xrightarrow{M_1}\ket{00}
         =\mathrm{SWAP}\ket{00}.\;\checkmark\\
\ket{01}&\xrightarrow{M_1}\ket{01}\xrightarrow{M_2}\ket{11}\xrightarrow{M_1}\ket{10}
         =\mathrm{SWAP}\ket{01}.\;\checkmark\\
\ket{10}&\xrightarrow{M_1}\ket{11}\xrightarrow{M_2}\ket{01}\xrightarrow{M_1}\ket{01}
         =\mathrm{SWAP}\ket{10}.\;\checkmark\\
\ket{11}&\xrightarrow{M_1}\ket{10}\xrightarrow{M_2}\ket{10}\xrightarrow{M_1}\ket{11}
         =\mathrm{SWAP}\ket{11}.\;\checkmark
\end{align*}
\end{solution}

\begin{exercice}[p.~31]
Is $\sqrt{1/3}\,(\ket{00}+\ket{01}+\ket{11})$ entangled?
And what about $\tfrac{1}{2}(\ket{00}+\ket{01}+\ket{10}+\ket{11})$?
\end{exercice}

\begin{solution}
\textbf{State 1: $\ket{\psi}=\frac{1}{\sqrt{3}}(\ket{00}+\ket{01}+\ket{11})$.}
Suppose $\ket{\psi}=(a\ket{0}+b\ket{1})\otimes(c\ket{0}+d\ket{1})=ac\ket{00}+ad\ket{01}+bc\ket{10}+bd\ket{11}$.
Matching coefficients: $ac=ad=bd=\tfrac{1}{\sqrt{3}}$ and $bc=0$.
From $bc=0$: either $b=0$ (then $bd=0$, contradiction) or $c=0$ (then $ac=0$, contradiction).
Hence $\ket{\psi}$ is \textbf{entangled}.

\medskip\noindent
\textbf{State 2: $\ket{\phi}=\tfrac{1}{2}(\ket{00}+\ket{01}+\ket{10}+\ket{11})$.}
We factor out the two qubits explicitly:
\begin{align*}
  \ket{\phi}
  &= \frac{1}{2}\bigl(\ket{00}+\ket{01}+\ket{10}+\ket{11}\bigr)\\
  &= \frac{1}{2}\bigl(\ket{0}(\ket{0}+\ket{1})+\ket{1}(\ket{0}+\ket{1})\bigr)\\
  &= \frac{1}{2}(\ket{0}+\ket{1})\otimes(\ket{0}+\ket{1})\\
  &= \frac{\ket{0}+\ket{1}}{\sqrt{2}}\otimes\frac{\ket{0}+\ket{1}}{\sqrt{2}}\\
  &= \ket{+}\otimes\ket{+}.
\end{align*}
Since $\ket{\phi}$ can be written as a tensor product of two single-qubit states,
it is a product state, hence \textbf{not entangled}.
\end{solution}

\begin{exercice}[p.~34]
Check that the matrix of $CU$ is
\[
  \begin{pmatrix}
    1 & 0 & 0           & 0           \\
    0 & 1 & 0           & 0           \\
    0 & 0 & u_{11}      & u_{12}      \\
    0 & 0 & u_{21}      & u_{22}
  \end{pmatrix}
\]
where $(u_{ij})_{i,j}$ is the matrix of $U$.
Check also that $CU$ is unitary. What is the adjoint of $CU$?
\end{exercice}

\begin{solution}
$CU$ acts as: $CU\ket{0}\ket{\psi}=\ket{0}\ket{\psi}$ (identity when control is $\ket{0}$) and
$CU\ket{1}\ket{\psi}=\ket{1}U\ket{\psi}$ (apply $U$ when control is $\ket{1}$).

In the ordered basis $\ket{00},\ket{01},\ket{10},\ket{11}$:
\begin{align*}
  CU\ket{00}&=\ket{00}
             =1\cdot\ket{00}+0\cdot\ket{01}+0\cdot\ket{10}+0\cdot\ket{11},\\
  CU\ket{01}&=\ket{01}
             =0\cdot\ket{00}+1\cdot\ket{01}+0\cdot\ket{10}+0\cdot\ket{11},\\
  CU\ket{10}&=\ket{1}U\ket{0}
             =\ket{1}(u_{11}\ket{0}+u_{21}\ket{1})
             =0\cdot\ket{00}+0\cdot\ket{01}+u_{11}\ket{10}+u_{21}\ket{11},\\
  CU\ket{11}&=\ket{1}U\ket{1}
             =\ket{1}(u_{12}\ket{0}+u_{22}\ket{1})
             =0\cdot\ket{00}+0\cdot\ket{01}+u_{12}\ket{10}+u_{22}\ket{11}.
\end{align*}
Each image gives a column of the matrix (the $k$-th column is the image of the $k$-th basis vector):
\[
  CU =
  \begin{pmatrix}
    1 & 0 & 0      & 0      \\
    0 & 1 & 0      & 0      \\
    0 & 0 & u_{11} & u_{12} \\
    0 & 0 & u_{21} & u_{22}
  \end{pmatrix}.\quad\checkmark
\]

\medskip\noindent
\textbf{Unitarity.}
The matrix is block-diagonal with blocks $I_2$ (top-left) and $U$ (bottom-right).
Computing $(CU)^\dagger(CU)$ explicitly:
\[
  (CU)^\dagger(CU)
  =\begin{pmatrix}I_2 & 0 \\ 0 & U^\dagger\end{pmatrix}
   \begin{pmatrix}I_2 & 0 \\ 0 & U\end{pmatrix}
  =\begin{pmatrix}I_2 I_2 & 0 \\ 0 & U^\dagger U\end{pmatrix}
  =\begin{pmatrix}I_2 & 0 \\ 0 & I_2\end{pmatrix}
  = I_4.\quad\checkmark
\]

\medskip\noindent
\textbf{Adjoint.}
\[
  (CU)^\dagger
  =\begin{pmatrix}I_2 & 0 \\ 0 & U^\dagger\end{pmatrix}
  = CU^\dagger.
\]
\end{solution}

\begin{exercice}[p.~35]
Prove the following equivalence: a controlled-$Z$ gate can be obtained from a
controlled-$X$ gate and two $H$ gates arranged as $H\,CX\,H$:
\[
\begin{quantikz}
  \qw & \ctrl{1} & \qw \\
  \qw & \gate{Z} & \qw
\end{quantikz}
\;=\;
\begin{quantikz}
  \qw & \ctrl{1} & \qw \\
  \gate{H} & \targ{} & \gate{H}
\end{quantikz}
\]
\end{exercice}

\begin{solution}

\textbf{Method 1 — matrices.}

\emph{Matrix of $CZ$.}
By definition, $CZ$ flips the sign of $\ket{11}$ and acts as identity on all other basis states:
\[
  CZ = \begin{pmatrix}1&0&0&0\\0&1&0&0\\0&0&1&0\\0&0&0&-1\end{pmatrix}.
\]

\emph{Matrix of $(I\otimes H)\,\mathrm{CNOT}\,(I\otimes H)$.}
We compute the two factors in turn.

The matrix of $I\otimes H$ (columns = images of basis vectors):
\[
  I\otimes H = \frac{1}{\sqrt{2}}\begin{pmatrix}1&1&0&0\\1&-1&0&0\\0&0&1&1\\0&0&1&-1\end{pmatrix}.
\]

The matrix of $\mathrm{CNOT}$ (control $=q_1$, target $=q_2$):
\[
  \mathrm{CNOT} = \begin{pmatrix}1&0&0&0\\0&1&0&0\\0&0&0&1\\0&0&1&0\end{pmatrix}.
\]

First compute $(I\otimes H)\cdot\mathrm{CNOT}$ by applying $I\otimes H$ to each column of $\mathrm{CNOT}$:
\begin{align*}
(I\otimes H)\,e_1 &= \tfrac{1}{\sqrt{2}}(1,1,0,0)^\top,\quad
(I\otimes H)\,e_2  = \tfrac{1}{\sqrt{2}}(1,-1,0,0)^\top,\\
(I\otimes H)\,e_4 &= \tfrac{1}{\sqrt{2}}(0,0,1,-1)^\top,\quad
(I\otimes H)\,e_3  = \tfrac{1}{\sqrt{2}}(0,0,1,1)^\top,
\end{align*}
giving
\[
  (I\otimes H)\,\mathrm{CNOT}
  = \frac{1}{\sqrt{2}}\begin{pmatrix}1&1&0&0\\1&-1&0&0\\0&0&1&1\\0&0&-1&1\end{pmatrix}.
\]

Then right-multiply by $I\otimes H$:
\[
  (I\otimes H)\,\mathrm{CNOT}\,(I\otimes H)
  = \frac{1}{2}
    \begin{pmatrix}1&1&0&0\\1&-1&0&0\\0&0&1&1\\0&0&-1&1\end{pmatrix}
    \begin{pmatrix}1&1&0&0\\1&-1&0&0\\0&0&1&1\\0&0&1&-1\end{pmatrix}.
\]
Computing row by row:
\begin{align*}
  \text{row }1:&\quad \tfrac{1}{2}(1{+}1,\,1{-}1,\,0,\,0) = (1,0,0,0),\\
  \text{row }2:&\quad \tfrac{1}{2}(1{-}1,\,1{+}1,\,0,\,0) = (0,1,0,0),\\
  \text{row }3:&\quad \tfrac{1}{2}(0,\,0,\,1{+}1,\,1{-}1) = (0,0,1,0),\\
  \text{row }4:&\quad \tfrac{1}{2}(0,\,0,\,-1{+}1,\,-1{-}1) = (0,0,0,-1).
\end{align*}
Hence
\[
  (I\otimes H)\,\mathrm{CNOT}\,(I\otimes H)
  = \begin{pmatrix}1&0&0&0\\0&1&0&0\\0&0&1&0\\0&0&0&-1\end{pmatrix}
  = CZ.\quad\checkmark
\]

\bigskip\noindent
\textbf{Method 2 — basis states.}

We check the action of $(I\otimes H)\,\mathrm{CNOT}\,(I\otimes H)$ on each basis state
and compare with $CZ$ (identity except $\ket{11}\mapsto-\ket{11}$).
\[
\ket{00}\xrightarrow{I\otimes H}\ket{0+}\xrightarrow{\mathrm{CNOT}}\ket{0+}\xrightarrow{I\otimes H}\ket{00}=CZ\ket{00}.\quad\checkmark
\]
\[
\ket{01}\xrightarrow{I\otimes H}\ket{0-}\xrightarrow{\mathrm{CNOT}}\ket{0-}\xrightarrow{I\otimes H}\ket{01}=CZ\ket{01}.\quad\checkmark
\]
\[
\ket{10}\xrightarrow{I\otimes H}\ket{1+}\xrightarrow{\mathrm{CNOT}}\ket{1}(X\ket{+})=\ket{1+}\xrightarrow{I\otimes H}\ket{10}=CZ\ket{10}.\quad\checkmark
\]
(Used $X\ket{+}=\ket{+}$ from Ex.~1.4.)
\[
\ket{11}\xrightarrow{I\otimes H}\ket{1-}\xrightarrow{\mathrm{CNOT}}-\ket{1-}\xrightarrow{I\otimes H}-\ket{11}=CZ\ket{11}.\quad\checkmark
\]
Since both circuits agree on every basis state, they are equal.\quad$\blacksquare$
\end{solution}

\begin{exercice}[p.~36]
Show that all four Bell states are entangled.
Obtain circuits to prepare them.
\textit{Hint}: you can use $Z$ and $X$ gates after the CNOT in the preceding
circuit.
\end{exercice}

\begin{solution}
The four Bell states are:
\[
  \ket{\Phi^\pm}=\frac{\ket{00}\pm\ket{11}}{\sqrt{2}},\qquad
  \ket{\Psi^\pm}=\frac{\ket{01}\pm\ket{10}}{\sqrt{2}}.
\]

\textbf{Entanglement.} Each state has the form $\sum_{a,b}\alpha_{ab}\ket{ab}$ with the coefficient matrix
$\begin{pmatrix}\alpha_{00}&\alpha_{01}\\\alpha_{10}&\alpha_{11}\end{pmatrix}$
having determinant $\alpha_{00}\alpha_{11}-\alpha_{01}\alpha_{10}\neq 0$
(it equals $\pm\tfrac{1}{2}$ for all four Bell states). A two-qubit state
$\ket{\psi}=(a\ket{0}+b\ket{1})\otimes(c\ket{0}+d\ket{1})$ would require this
determinant to be $acbd-adbc=0$. Hence none of the Bell states is a product state:
all four are \textbf{entangled}.

\medskip\noindent
\textbf{Preparation circuits.}
Starting from $\ket{00}$, the base circuit $(H\otimes I)$ then $\mathrm{CNOT}$ gives $\ket{\Phi^+}$.
The three remaining Bell states follow by appending correction gates \emph{after} the CNOT:

\[
\begin{array}{cl}
\ket{\Phi^+}: &
  \ket{00}\xrightarrow{H\otimes I}\dfrac{1}{\sqrt{2}}(\ket{00}+\ket{10})
  \xrightarrow{\mathrm{CNOT}}\dfrac{1}{\sqrt{2}}(\ket{00}+\ket{11}).\\[8pt]
\ket{\Phi^-}: &
  \dfrac{1}{\sqrt{2}}(\ket{00}+\ket{11})
  \xrightarrow{\mathrm{CZ}}\dfrac{1}{\sqrt{2}}(\ket{00}-\ket{11}).\\[8pt]
\ket{\Psi^+}: &
  \dfrac{1}{\sqrt{2}}(\ket{00}+\ket{11})
  \xrightarrow{X\otimes I}\dfrac{1}{\sqrt{2}}(\ket{10}+\ket{01})
  =\dfrac{1}{\sqrt{2}}(\ket{01}+\ket{10}).\\[8pt]
\ket{\Psi^-}: &
  \dfrac{1}{\sqrt{2}}(\ket{00}-\ket{11})
  \xrightarrow{X\otimes I}\dfrac{1}{\sqrt{2}}(\ket{10}-\ket{01})
  =-\ket{\Psi^-}.
\end{array}
\]

\noindent The last result equals $\ket{\Psi^-}$ up to a global phase of $-1$.
In summary, writing gates in order of application from $\ket{00}$:
\[
\begin{array}{ll}
\ket{\Phi^+}: & (H\otimes I),\ \mathrm{CNOT}.\\[3pt]
\ket{\Phi^-}: & (H\otimes I),\ \mathrm{CNOT},\ \mathrm{CZ}.\\[3pt]
\ket{\Psi^+}: & (H\otimes I),\ \mathrm{CNOT},\ (X\otimes I).\\[3pt]
\ket{\Psi^-}: & (H\otimes I),\ \mathrm{CNOT},\ \mathrm{CZ},\ (X\otimes I).
\end{array}
\]
\end{solution}

\begin{exercice}[p.~38]
Check that the basis state $\ket{j}$ is represented by a $2^n$-dimensional
column vector whose $j$-th component is $1$, while the rest are $0$
(\textit{Hint}: Use, repeatedly, the expression for the tensor product of
column vectors that we discussed in \textit{Section~1.4.1} and the fact that
the tensor product is associative).
Deduce that any $n$-qubit state can be represented by a $2^n$-dimensional
column vector with unit length.
\end{exercice}

\begin{solution}
\textbf{Base case} ($n=1$): By definition,
\[
  \ket{0}=\begin{pmatrix}1\\0\end{pmatrix},\qquad
  \ket{1}=\begin{pmatrix}0\\1\end{pmatrix}.
\]
For $j=0$: the component at index $0$ is $1$, the component at index $1$ is $0$.\quad$\checkmark$\\
For $j=1$: the component at index $1$ is $1$, the component at index $0$ is $0$.\quad$\checkmark$

\medskip\noindent
\textbf{Inductive step}: Assume the claim holds for $n-1$ qubits: for every
$j\in\{0,\ldots,2^{n-1}-1\}$, the basis state $\ket{j}$ is the $2^{n-1}$-dimensional
vector with a $1$ at position $j$ and $0$ everywhere else.

Every $n$-qubit basis state has the form $\ket{j_1}\otimes\ket{j}$
with $j_1\in\{0,1\}$ and $j\in\{0,\ldots,2^{n-1}-1\}$.
Applying the rule $\binom{a}{b}\otimes v=\binom{av}{bv}$:

\medskip
\textit{Case $j_1=0$:}
\[
  \ket{0}\otimes\ket{j}
  =\begin{pmatrix}1\\0\end{pmatrix}\otimes\ket{j}
  =\begin{pmatrix}\ket{j}\\[2pt]\mathbf{0}\end{pmatrix}.
\]
The top half of the $2^n$-vector is $\ket{j}$, the bottom half is all zeros.
By the IH, $\ket{j}$ has its $1$ at position $j$, so the full vector has its $1$ at position $j$.\quad$\checkmark$

\medskip
\textit{Case $j_1=1$:}
\[
  \ket{1}\otimes\ket{j}
  =\begin{pmatrix}0\\1\end{pmatrix}\otimes\ket{j}
  =\begin{pmatrix}\mathbf{0}\\[2pt]\ket{j}\end{pmatrix}.
\]
The top half is all zeros, the bottom half is $\ket{j}$.
By the IH, $\ket{j}$ has its $1$ at position $j$ within the bottom half,
so the full vector has its $1$ at position $2^{n-1}+j$.\quad$\checkmark$

\medskip\noindent
Ranging over both cases and all $j\in\{0,\ldots,2^{n-1}-1\}$, the indices $j$ (case~0)
and $2^{n-1}+j$ (case~1) together cover $\{0,\ldots,2^n-1\}$ without overlap:
every $n$-qubit basis state is a distinct standard basis vector of $\mathbb{C}^{2^n}$.\quad$\checkmark$

\medskip\noindent
\textbf{Example} ($n=2$): using $\binom{a}{b}\otimes\binom{c}{d}=\bigl(ac,\,ad,\,bc,\,bd\bigr)^T$,
\[
  \ket{00}=\begin{pmatrix}1\\0\\0\\0\end{pmatrix},\quad
  \ket{01}=\begin{pmatrix}0\\1\\0\\0\end{pmatrix},\quad
  \ket{10}=\begin{pmatrix}0\\0\\1\\0\end{pmatrix},\quad
  \ket{11}=\begin{pmatrix}0\\0\\0\\1\end{pmatrix}.
\]
For $j_1=0$: the $1$ is in the top half at position $j_2\in\{0,1\}$.
For $j_1=1$: the $1$ is in the bottom half at position $2+j_2\in\{2,3\}$.\quad$\checkmark$

\medskip\noindent
\textbf{Orthonormality of the basis.}
Since each $\ket{j}=e_j$ has a $1$ at position $j$ and $0$ elsewhere,
the inner product of two basis vectors is
\[
  \braket{j|k} = e_j^\dagger\, e_k = \delta_{jk},
\]
because the only non-zero entry of $e_j$ is at position $j$: the dot product is $1$ if $j=k$ and $0$ otherwise.
Hence the $2^n$ basis states form an \textbf{orthonormal} family.\quad$\checkmark$

\medskip\noindent
\textbf{Deduction}: Any $n$-qubit state can be expanded in this orthonormal basis as
$\ket{\psi}=\sum_{j=0}^{2^n-1}\alpha_j\ket{j}$, where $\alpha_j=\braket{j|\psi}$.
Since $\ket{j}=e_j$, linearity gives
\[
  \ket{\psi}=\sum_{j=0}^{2^n-1}\alpha_j\,e_j
  =\begin{pmatrix}\alpha_0\\\alpha_1\\\vdots\\\alpha_{2^n-1}\end{pmatrix}\in\mathbb{C}^{2^n}.
\]
The normalisation condition $\braket{\psi|\psi}=1$ reads
\[
  \sum_{j,k}\alpha_j^*\alpha_k\,\braket{j|k}
  =\sum_{j,k}\alpha_j^*\alpha_k\,\delta_{jk}
  =\sum_j|\alpha_j|^2=1,
\]
which is exactly the unit $\ell^2$-norm condition for this column vector.\quad$\checkmark$
\end{solution}

\begin{exercice}[p.~38]
Consider an $n$-qubit state $\ket{\psi}=\sum_{j\in\{0,1\}^n}\alpha_j\ket{j}$.
Define $J_0=\{j\in\{0,1\}^n : j_1=0\}$ as the set of basis indices whose first bit is $0$.
When the first qubit is measured and the outcome is $0$, the book gives:
\[
  P(\text{first qubit}=0)=\sum_{j\in J_0}|\alpha_j|^2,
  \qquad
  \ket{\psi'_0}=\frac{\displaystyle\sum_{j\in J_0}\alpha_j\ket{j}}
                     {\sqrt{P(\text{first qubit}=0)}}.
\]
Derive the analogous formulas for the case in which the outcome is $1$,
using the set $J_1=\{j\in\{0,1\}^n : j_1=1\}$.
\end{exercice}

\begin{solution}
By symmetry with the outcome-$0$ case, replacing $J_0$ by $J_1=\{j\in\{0,1\}^n:j_1=1\}$:
\[
  P(\text{first qubit}=1)=\sum_{j\in J_1}|\alpha_j|^2,
  \qquad
  \ket{\psi'_1}=\frac{\displaystyle\sum_{j\in J_1}\alpha_j\ket{j}}
                     {\sqrt{P(\text{first qubit}=1)}}.
\]
Note that $J_0$ and $J_1$ partition $\{0,1\}^n$, so
$P(\text{first qubit}=0)+P(\text{first qubit}=1)=\sum_{j}|\alpha_j|^2=1$.\quad$\checkmark$
\end{solution}

\begin{exercice}[p.~38]
What is the probability of getting $0$ when we measure the second qubit of
$(1/2)\ket{100} + (1/2)\ket{010} + \sqrt{1/2}\,\ket{001}$?
What will the state be after the measurement if we indeed get $0$?
\end{exercice}

\begin{solution}
Write the state $\ket{\psi}=\tfrac{1}{2}\ket{100}+\tfrac{1}{2}\ket{010}+\tfrac{1}{\sqrt{2}}\ket{001}$.
The second qubit (middle bit) equals $0$ in $\ket{100}$ and $\ket{001}$, and equals $1$ in $\ket{010}$.

\[
  P(q_2=0) = \left|\tfrac{1}{2}\right|^2+\left|\tfrac{1}{\sqrt{2}}\right|^2
           = \frac{1}{4}+\frac{1}{2} = \frac{3}{4}.
\]
Post-measurement state (keeping only the terms with $q_2=0$ and renormalising):
\[
  \ket{\psi'} = \frac{\frac{1}{2}\ket{100}+\frac{1}{\sqrt{2}}\ket{001}}{\sqrt{3/4}}
              = \frac{2}{\sqrt{3}}\left(\frac{1}{2}\ket{100}+\frac{1}{\sqrt{2}}\ket{001}\right)
              = \frac{1}{\sqrt{3}}\ket{100}+\sqrt{\frac{2}{3}}\ket{001}.
\]
Check: $\tfrac{1}{3}+\tfrac{2}{3}=1$.\quad$\checkmark$
\end{solution}

\begin{exercice}[p.~39]
Compute the inner product of $\ket{x}$ and $\ket{y}$, where $x$ and $y$ are
both binary strings of length $n$. Use your result to prove that
$\bigl\{\ket{x}\bigr\}_{x\in\{0,1\}^n}$ is, indeed, an orthonormal basis.
\end{exercice}

\begin{solution}
Write $\ket{x}=\sum_{i}x_i\ket{i}$ and $\ket{y}=\sum_{j}y_j\ket{j}$
(sums over $i,j\in\{0,\ldots,2^n-1\}$),
where the components satisfy $x_i=1$ if $i=x$ and $0$ otherwise (and likewise for $y_j$).

Since $y_j\in\{0,1\}\subset\mathbb{R}$, we have $\bar{y}_j=y_j$.
Using $\braket{j|i}=\delta_{i,j}$:
\[
  \braket{y|x}
  =\sum_{j}\sum_{i}y_j\,x_i\,\braket{j|i}
  =\sum_{j}\sum_{i}y_j\,x_i\,\delta_{i,j}
  =\sum_{i}y_i\,x_i.
\]

\textit{Case $x=y$}: the only nonzero term in $\sum_i y_i x_i$ is $i=x=y$,
contributing $y_x\cdot x_x=1\cdot 1=1$.
Hence $\braket{x|x}=1$.\quad$\checkmark$

\textit{Case $x\neq y$}: $x_i=1$ only at $i=x$, and $y_i=1$ only at $i=y$.
Since $x\neq y$ these indices never coincide, so every term $y_i x_i=0$.
Hence $\braket{y|x}=0$.\quad$\checkmark$

In both cases, $\braket{y|x}=\delta_{x,y}$: the $2^n$ basis states are \textbf{orthonormal}.
Since there are $2^n$ of them in a $2^n$-dimensional space, they form a \textbf{basis}.\quad$\blacksquare$
\end{solution}

\begin{exercice}[p.~39]
Compute the inner product of the states
$\sqrt{1/2}\,(\ket{000}+\ket{111})$ and
$\tfrac{1}{2}(\ket{000}+\ket{011}+\ket{101}+\ket{110})$.
\end{exercice}

\begin{solution}
Let $\ket{\psi_1}=\tfrac{1}{\sqrt{2}}(\ket{000}+\ket{111})$ and
$\ket{\psi_2}=\tfrac{1}{2}(\ket{000}+\ket{011}+\ket{101}+\ket{110})$.

Using $\braket{x|y}=\delta_{x,y}$:
\[
  \braket{\psi_1|\psi_2}
  =\frac{1}{\sqrt{2}}\cdot\frac{1}{2}\bigl(\braket{000|000}+\underbrace{\braket{000|011}}_{0}
    +\cdots+\underbrace{\braket{111|110}}_{0}\bigr)
  =\frac{1}{2\sqrt{2}}\cdot 1 = \frac{1}{2\sqrt{2}}=\frac{\sqrt{2}}{4}.
\]
Only the $\braket{000|000}$ term contributes (the string $111$ does not appear in $\ket{\psi_2}$).
\end{solution}

\begin{exercice}[p.~40]
Obtain the matrix for the CCNOT gate and verify that it is unitary.
\end{exercice}

\begin{solution}
\textbf{Definition.}
The Toffoli gate acts as
\[
  \mathrm{CCNOT}\ket{q_1\,q_2\,q_3}=\ket{q_1\;q_2\;(q_3\oplus(q_1\wedge q_2))}.
\]
It flips $q_3$ if and only if $q_1=q_2=1$; every other basis state is unchanged.

\medskip\noindent
\textbf{Action on all basis states:}
\[
\begin{array}{r@{\;\mapsto\;}l@{\qquad\qquad}r@{\;\mapsto\;}l}
\ket{000} & \ket{000} & \ket{100} & \ket{100} \\[2pt]
\ket{001} & \ket{001} & \ket{101} & \ket{101} \\[2pt]
\ket{010} & \ket{010} & \ket{110} & \ket{111} \\[2pt]
\ket{011} & \ket{011} & \ket{111} & \ket{110}
\end{array}
\]

\medskip\noindent
\textbf{Deriving the matrix.}
The $j$-th column of the matrix is the vector $\mathrm{CCNOT}\ket{j}$.
Each output above is itself a basis state, so each column is a standard basis vector $e_k$
(with $1$ at row $k$ and $0$ elsewhere).
Reading off all eight columns:

\[
\begin{array}{c|l}
\text{column} & \text{vector}\\
\hline
0\ (\ket{000}\mapsto\ket{000}) & e_0=(1,0,0,0,0,0,0,0)^T \\
1\ (\ket{001}\mapsto\ket{001}) & e_1=(0,1,0,0,0,0,0,0)^T \\
2\ (\ket{010}\mapsto\ket{010}) & e_2=(0,0,1,0,0,0,0,0)^T \\
3\ (\ket{011}\mapsto\ket{011}) & e_3=(0,0,0,1,0,0,0,0)^T \\
4\ (\ket{100}\mapsto\ket{100}) & e_4=(0,0,0,0,1,0,0,0)^T \\
5\ (\ket{101}\mapsto\ket{101}) & e_5=(0,0,0,0,0,1,0,0)^T \\
6\ (\ket{110}\mapsto\ket{111}) & e_7=(0,0,0,0,0,0,0,1)^T \\
7\ (\ket{111}\mapsto\ket{110}) & e_6=(0,0,0,0,0,0,1,0)^T
\end{array}
\]

Assembling these as columns:
\[
  \mathrm{CCNOT}=
  \begin{pmatrix}
    1&0&0&0&0&0&0&0\\
    0&1&0&0&0&0&0&0\\
    0&0&1&0&0&0&0&0\\
    0&0&0&1&0&0&0&0\\
    0&0&0&0&1&0&0&0\\
    0&0&0&0&0&1&0&0\\
    0&0&0&0&0&0&0&1\\
    0&0&0&0&0&0&1&0
  \end{pmatrix}
  =\begin{pmatrix}I_6 & 0 \\ 0 & X\end{pmatrix}.
\]

\medskip\noindent
\textbf{Unitarity.}
The matrix is block-diagonal with $I_6$ in the top-left block and
$X=\bigl(\begin{smallmatrix}0&1\\1&0\end{smallmatrix}\bigr)$ in the bottom-right block.
Both blocks are unitary ($I_6^\dagger I_6=I_6$ and $X^\dagger X=I_2$),
so $\mathrm{CCNOT}^\dagger\mathrm{CCNOT}=I_8$.\quad$\checkmark$

Alternatively, CCNOT is its own inverse: applying it twice gives
$q_3\oplus(q_1\wedge q_2)\oplus(q_1\wedge q_2)=q_3$,
so $\mathrm{CCNOT}^2=I$, hence $\mathrm{CCNOT}^{-1}=\mathrm{CCNOT}=\mathrm{CCNOT}^\dagger$.\quad$\checkmark$
\end{solution}

\begin{exercice}[p.~41]
Verify that the following circuit implements the Toffoli gate by checking
its action on the states of the computational basis.

% Right side: Decomposition (Mirrored vertically)
\begin{quantikz}
\gate{H} & \targ{} & \gate{T^\dagger} & \targ{} & \gate{T} & \targ{} & \gate{T^\dagger} & \targ{} & \gate{T} & \gate{H} & \qw & \qw & \qw \\
\qw      & \ctrl{-1} & \qw             & \qw      & \qw      & \ctrl{-1} & \qw             & \qw      & \gate{T} & \targ{} & \gate{T^\dagger} & \targ{} & \qw \\
\qw      & \qw      & \qw             & \ctrl{-2} & \qw      & \qw      & \qw             & \ctrl{-2} & \qw      & \ctrl{-1} & \gate{T} & \ctrl{-1} & \qw
\end{quantikz}
\end{exercice}

\begin{adjustwidth}{-1.5cm}{-1.5cm}
\begin{solution}
The circuit uses qubits in order $(z,y,x)$ where $z$ is the top wire.
It implements Toffoli: $z$ is the target, $y$ and $x$ are the controls (flip $z$ iff $x=y=1$).

\medskip
\noindent\textbf{State evolution:}

\medskip
\centering\resizebox{\linewidth}{!}{%
\begin{footnotesize}%
\begin{tabular}{@{}r@{\;}l@{\quad}l@{\quad}l@{\quad}l@{}}
\toprule
 & \textbf{Operator} & $\ket{000}$ (idle) & $\ket{110}$ & $\ket{111}$ \\
\midrule
0  & (input)
   & $\ket{000}$
   & $\ket{110}$ & $\ket{111}$ \\[3pt]
1  & $H \otimes I \otimes I$
   & $\tfrac{1}{\sqrt{2}}(\ket{000}+\ket{100})$
   & $\tfrac{1}{\sqrt{2}}(\ket{010}-\ket{110})$
   & $\tfrac{1}{\sqrt{2}}(\ket{011}-\ket{111})$ \\[3pt]
2  & $(|0\rangle\langle0|\!\otimes\! I\!\otimes\! I)+(|1\rangle\langle1|\!\otimes\! X\!\otimes\! I)$
   & $\tfrac{1}{\sqrt{2}}(\ket{000}+\ket{110})$
   & $\tfrac{1}{\sqrt{2}}(\ket{010}-\ket{100})$
   & $\tfrac{1}{\sqrt{2}}(\ket{011}-\ket{101})$ \\[3pt]
3  & $T^\dagger \otimes I \otimes I$
   & $\tfrac{1}{\sqrt{2}}(\ket{000}+e^{-i\pi/4}\ket{110})$
   & $\tfrac{1}{\sqrt{2}}(\ket{010}-e^{-i\pi/4}\ket{100})$
   & $\tfrac{1}{\sqrt{2}}(\ket{011}-e^{-i\pi/4}\ket{101})$ \\[3pt]
4  & $(|0\rangle\langle0|\!\otimes\! I\!\otimes\! I)+(|1\rangle\langle1|\!\otimes\! I\!\otimes\! X)$
   & $\tfrac{1}{\sqrt{2}}(\ket{000}+e^{-i\pi/4}\ket{111})$
   & $\tfrac{1}{\sqrt{2}}(\ket{010}-e^{-i\pi/4}\ket{101})$
   & $\tfrac{1}{\sqrt{2}}(\ket{011}-e^{-i\pi/4}\ket{100})$ \\[3pt]
5  & $T \otimes I \otimes I$
   & $\tfrac{1}{\sqrt{2}}(\ket{000}+\ket{111})$
   & $\tfrac{1}{\sqrt{2}}(\ket{010}-\ket{101})$
   & $\tfrac{1}{\sqrt{2}}(\ket{011}-\ket{100})$ \\[3pt]
6  & $(|0\rangle\langle0|\!\otimes\! I\!\otimes\! I)+(|1\rangle\langle1|\!\otimes\! X\!\otimes\! I)$
   & $\tfrac{1}{\sqrt{2}}(\ket{000}+\ket{101})$
   & $\tfrac{1}{\sqrt{2}}(\ket{010}-\ket{111})$
   & $\tfrac{1}{\sqrt{2}}(\ket{011}-\ket{110})$ \\[3pt]
7  & $T^\dagger \otimes I \otimes I$
   & $\tfrac{1}{\sqrt{2}}(\ket{000}+e^{-i\pi/4}\ket{101})$
   & $\tfrac{1}{\sqrt{2}}(\ket{010}-e^{-i\pi/4}\ket{111})$
   & $\tfrac{1}{\sqrt{2}}(\ket{011}-e^{-i\pi/4}\ket{110})$ \\[3pt]
8  & $(|0\rangle\langle0|\!\otimes\! I\!\otimes\! I)+(|1\rangle\langle1|\!\otimes\! I\!\otimes\! X)$
   & $\tfrac{1}{\sqrt{2}}(\ket{000}+e^{-i\pi/4}\ket{100})$
   & $\tfrac{1}{\sqrt{2}}(\ket{010}-e^{-i\pi/4}\ket{110})$
   & $\tfrac{1}{\sqrt{2}}(\ket{011}-e^{-i\pi/4}\ket{111})$ \\[3pt]
9  & $T \otimes T \otimes I$
   & $\tfrac{1}{\sqrt{2}}(\ket{000}+\ket{100})$
   & $\tfrac{1}{\sqrt{2}}(\ket{010}-e^{i\pi/4}\ket{110})$
   & $\tfrac{1}{\sqrt{2}}(\ket{011}-e^{i\pi/4}\ket{111})$ \\[3pt]
10 & $H\!\otimes\!(|0\rangle\langle0|\!\otimes\! I+|1\rangle\langle1|\!\otimes\! X)$
   & $\ket{000}$
   & $\ket{110}$
   & $e^{i\pi/4}\tfrac{1}{\sqrt{2}}(\ket{010}+\ket{110})$ \\[3pt]
11 & $I \otimes T^\dagger \otimes T$
   & $\ket{000}$
   & $\ket{110}$
   & $\tfrac{1}{\sqrt{2}}(\ket{010}+e^{i\pi/4}\ket{110})$ \\[3pt]
12 & $I\!\otimes\!(|0\rangle\langle0|\!\otimes\! I+|1\rangle\langle1|\!\otimes\! X)$
   & $\ket{000}$
   & $\ket{110}$
   & $\ket{011}$ \\
\bottomrule
\end{tabular}%
\end{footnotesize}}%



\medskip\noindent
\textbf{Cases $q_1=0$ or $q_2=0$ (no Toffoli flip should occur).}
When the first control $q_1=0$, the gates $\mathrm{CNOT}_{q_1\to q_3}$ are idle.
The sub-circuit acting on $q_3$ then contains $H,T^\dagger,T,T^\dagger,T,H$ with two
$\mathrm{CNOT}_{q_2\to q_3}$ interspersed, but the paired $T^\dagger/T$ cancellations yield the identity.
A careful term-by-term computation (or observing that the circuit reduces to
$H(T^\dagger T)(T^\dagger T)H=HH=I$ when the CNOTs are idle) gives the identity on $q_3$.
Similarly when $q_2=0$.  Hence the eight cases with $q_1=0$ or $q_2=0$ all return the
input state unchanged, consistent with the Toffoli gate.\quad$\checkmark$

\medskip\noindent
\textbf{Case $\ket{110}$ ($q_1=q_2=1$, $q_3=0$).}
We trace $\ket{q_3}=\ket{0}$ through the 10 gates acting on $q_3$ (with both controls active):
\begin{align*}
  \ket{0}&\xrightarrow{H}\tfrac{1}{\sqrt{2}}(\ket{0}+\ket{1})
          \xrightarrow{T^\dagger}\tfrac{1}{\sqrt{2}}(\ket{0}+c^*\ket{1})
          \xrightarrow{T\text{ (CNOT flips)}}\tfrac{1}{\sqrt{2}}(\ket{0}+c\cdot c^*\ket{1})
          =\tfrac{1}{\sqrt{2}}(\ket{0}+\ket{1})\\
         &\xrightarrow{T^\dagger}\tfrac{1}{\sqrt{2}}(\ket{0}+c^*\ket{1})
          \xrightarrow{T\text{ (CNOT flips)}}\tfrac{1}{\sqrt{2}}(\ket{0}+\ket{1})
          \xrightarrow{T}\tfrac{1}{\sqrt{2}}(\ket{0}+c\ket{1})
          \xrightarrow{H}\cdots=\ket{1}.
\end{align*}
(The two CNOT$_{q_2\to q_3}$ and two CNOT$_{q_1\to q_3}$ contribute phase factors that combine
to advance the phase by $c^4=-1$ relative to the $\ket{0}$ amplitude, giving $\ket{1}$ after the
final $H$.)  The last CNOT$_{q_1\to q_2}$ section contributes identity on $q_2$ (since $c\cdot c^*=1$).
Result: $\ket{110}\to\ket{111}=$ Toffoli$\ket{110}$.\quad$\checkmark$

\medskip\noindent
\textbf{Case $\ket{111}$ ($q_1=q_2=1$, $q_3=1$).}
An identical trace starting from $\ket{1}$ yields $\ket{0}$ after the target-qubit sub-circuit
(the phase accumulation is the same but now starts on $\ket{1}$, reversing the output).
Result: $\ket{111}\to\ket{110}=$ Toffoli$\ket{111}$.\quad$\checkmark$

\medskip\noindent
The circuit therefore acts as the Toffoli gate on all $8$ basis states.\quad$\blacksquare$
\end{solution}
\end{adjustwidth}


%%%%%%%%%%%%%%%%%%%%%%%%%%%%%%%%%%%%%%%%%%%%%%%%%%%
%% Fin du document
%%%%%%%%%%%%%%%%%%%%%%%%%%%%%%%%%%%%%%%%%%%%%%%%%%%
\end{document}
